\section{Results}
\label{sec:results}

\subsection{Cohort characteristics and descriptive outcomes}

The analytic cohort included \CohortN{} ED encounters; \CohortEvents{} (56.3\%) had at least one documented pain reassessment during the stay (Table~\ref{tab:cohort}).
The mean age was 56.2 years (SD 19.7); 70.4\% of patients were White, 19.7\% Black, 6.9\% Hispanic, and 3.1\% Asian.
Most encounters were trauma-related (92.8\%); acute pancreatitis accounted for 7.2\%.
Insurance distribution included Medicare, private, Medicaid, undocumented (missing payer in EHR), and uninsured (self-pay); see Table~\ref{tab:cohort}.
Mean initial pain score was 6.1 (SD 2.6); mean ESI was 2.7 (SD 0.7).

Reassessment was often delayed: the median time to first reassessment was 185 minutes (IQR 96--311).
Only 9.6\% of patients were reassessed within 60 minutes and 21.6\% within 120 minutes.
Unadjusted Kaplan--Meier curves showed separation by insurance and race/ethnicity (Figure~\ref{fig:km}); 60-minute reassessment rates varied by pain severity stratum (Figure~\ref{fig:rates60}).

\begin{figure}[htbp]
  \centering
  \includegraphics[width=0.85\linewidth]{fig02_cohort_flow.png}
  \caption{Cohort construction for the primary analytic sample.}
  \label{fig:flow}
\end{figure}

\begin{table}[htbp]
  \centering
  \caption{Characteristics of the analytic cohort ($N=\CohortN{}$)}
  \label{tab:cohort}
  \footnotesize
  \begin{tabular}{@{}lc@{}}
    \toprule
    Characteristic & Value \\
    \midrule
    \multicolumn{2}{@{}l}{\textit{Demographics}} \\
    Age (years), mean (SD) & 56.2 (19.7) \\
    \quad Age 18--39, $n$ (\%) & 3,897 (22.7) \\
    \quad Age 40--64, $n$ (\%) & 7,057 (41.0) \\
    \quad Age 65+, $n$ (\%) & 6,250 (36.3) \\
    Female sex, $n$ (\%) & 9,797 (56.9) \\
    White race/ethnicity, $n$ (\%) & 12,108 (70.4) \\
    Black race/ethnicity, $n$ (\%) & 3,383 (19.7) \\
    Hispanic ethnicity, $n$ (\%) & 1,186 (6.9) \\
    Asian race/ethnicity, $n$ (\%) & 527 (3.1) \\
    \midrule
    \multicolumn{2}{@{}l}{\textit{Insurance}} \\
    Medicare, $n$ (\%) & 8,015 (46.6) \\
    Private, $n$ (\%) & 5,071 (29.5) \\
    Medicaid, $n$ (\%) & 4,118 (23.9) \\
    \midrule
    \multicolumn{2}{@{}l}{\textit{Clinical}} \\
    ESI triage level, mean (SD) & 2.7 (0.7) \\
    Initial pain score, mean (SD) & 6.1 (2.6) \\
    Trauma diagnosis, $n$ (\%) & 15,958 (92.8) \\
    Acute pancreatitis, $n$ (\%) & 1,246 (7.2) \\
    \midrule
    \multicolumn{2}{@{}l}{\textit{Outcomes}} \\
    Any reassessment, $n$ (\%) & 10,569 (61.4) \\
    Reassessed within 60 min, $n$ (\%) & 1,651 (9.6) \\
    Reassessed within 120 min, $n$ (\%) & 3,722 (21.6) \\
    Minutes to reassessment, median (IQR) & 185 (96--311) \\
  \bottomrule
  \end{tabular}
\end{table}


\begin{figure}[htbp]
  \centering
  \includegraphics[width=\linewidth]{fig04_km_reassessment_overview.png}
  \caption{Unadjusted cumulative probability of first pain reassessment by race/ethnicity, insurance, and ESI.}
  \label{fig:km}
\end{figure}

\begin{figure}[htbp]
  \centering
  \includegraphics[width=0.9\linewidth]{fig05_reassessment_rates_60min.png}
  \caption{Unadjusted proportion reassessed within 60 minutes by initial pain severity group and race/ethnicity or insurance.}
  \label{fig:rates60}
\end{figure}

\subsection{Primary adjusted model: clinical severity and workflow}

In the primary M4 Cox model (Figure~\ref{fig:m4forest}), clinical presentation and ED workflow factors were strongly associated with reassessment timing.
Each one-point increase in initial pain score was associated with faster reassessment (HR \HRPainMfour{}).
Higher ESI level (lower acuity) was associated with slower reassessment (HR \HRESIMfour{}).
Ambulance arrival (vs.\ walk-in) and night shift (vs.\ day) were associated with faster reassessment (HR \HRAmbulanceMfour{} and HR \HRNightMfour{}, respectively).
Weekend arrival and ED crowding variables showed weaker or null associations at conventional significance thresholds.

\begin{figure}[htbp]
  \centering
  \includegraphics[width=0.95\linewidth]{fig06_m4_sectional_forest.png}
  \caption{Primary adjusted Cox model (M4): hazard ratios for time to first pain reassessment.
  HR $>1$: faster reassessment; HR $<1$: slower reassessment.
  Vitals are adjusted in M3 but not displayed; other arrival mode is excluded from the main figure (see supplement).}
  \label{fig:m4forest}
\end{figure}

\subsection{Insurance: strongest and most persistent disparity signal}

Insurance showed the largest and most robust associations with reassessment timing.
Compared with private insurance, Medicaid was associated with substantially slower reassessment in M4 (HR \HRMedicaidMfour{}).
Medicare was also associated with slower reassessment (HR \HRMedicareMfour{}).
Undocumented insurance (missing payer documentation) had a hazard ratio near unity versus private (HR \HRUndocumentedMfour{}).

Sequential attenuation demonstrated minimal change in the Medicaid association from M2 (\HRMedicaidMtwo{}) through M3 (\HRMedicaidMthree{}) to M4 (\HRMedicaidMfourSeq{}; Figure~\ref{fig:attenuation}).
An insurance-focused analysis including Kaplan--Meier curves, within-acuity estimates, and pain-stratum models corroborated this pattern (Figure~\ref{fig:insurance}).

\begin{figure}[htbp]
  \centering
  \includegraphics[width=\linewidth]{fig10_insurance_focused_analysis.png}
  \caption{Insurance-focused analysis: unadjusted curves, M4 hazard ratios, within-acuity Medicaid estimates, and pain-stratum models.}
  \label{fig:insurance}
\end{figure}

\subsection{Race and ethnicity: smaller effects with attenuation}

Race and ethnicity associations were generally smaller than insurance and attenuated with sequential adjustment.
Black patients had slower reassessment in M2 (HR \HRBlackMtwo{}) that moved toward the null in M4 (HR \HRBlackMfour{}; Figure~\ref{fig:attenuation}).
Hispanic patients showed no meaningful association after full adjustment (HR \HRHispanicMfour{}).
Asian patients had slower reassessment in M4 (HR \HRAsianMfour{}), but estimates were less stable in severe-pain sensitivity analyses (wide confidence intervals).

\begin{figure}[htbp]
  \centering
  \includegraphics[width=\linewidth]{fig08_sequential_attenuation_key_factors.png}
  \caption{Sequential attenuation of key factors across models M1--M4 (race, insurance, initial pain).}
  \label{fig:attenuation}
\end{figure}

\subsection{Within-acuity and severe-pain sensitivity analyses}

Within-acuity models (Figure~\ref{fig:withinacuity}) suggested that insurance associations were most pronounced at higher acuity.
Among ESI 1--2 encounters, Medicaid was associated with slower reassessment (HR \HRMedicaidESI12{}).
Among ESI 3 encounters, Medicaid remained associated with slower reassessment (HR \HRMedicaidESI3{}).
Estimates in ESI 4--5 were imprecise and directionally unstable, reflecting smaller event counts.

In severe-pain sensitivity analyses (initial pain 7--10; Figure~\ref{fig:severe}), patterns were similar to the primary cohort: Medicaid HR \HRMedicaidSevere{} and Black vs.\ White HR \HRBlackSevere{}.

\begin{figure}[htbp]
  \centering
  \includegraphics[width=\linewidth]{fig11_within_acuity_forests.png}
  \caption{Within-acuity M4-style models: race/ethnicity and insurance by ESI stratum.}
  \label{fig:withinacuity}
\end{figure}

\begin{figure}[htbp]
  \centering
  \includegraphics[width=0.95\linewidth]{fig12_severe_pain_sensitivity.png}
  \caption{Severe-pain sensitivity analysis (pain 7--10): key M4 hazard ratios.}
  \label{fig:severe}
\end{figure}

\subsection{Temporal trends}

We observed no evidence of a meaningful secular trend in reassessment timing.
In a supplementary Cox model with continuous de-identified year (M4 covariates), the HR per 5-year increase was \HRYearFive{}.
Figure~\ref{fig:yeartrend} displays unadjusted and M4-adjusted probabilities of 60-minute reassessment by policy era; both series fluctuated between approximately 7\% and 12\% without a consistent directional pattern.

\begin{figure}[htbp]
  \centering
  \includegraphics[width=0.95\linewidth]{fig07_year_era_reassessment_trend.png}
  \caption{Temporal trend in 60-minute reassessment by de-identified policy era: unadjusted observed rates (bars) and M4-adjusted predicted probabilities (points) with 95\% confidence intervals.}
  \label{fig:yeartrend}
\end{figure}

\subsection{Post-analgesic pathway analysis}

In post-analgesic sensitivity analyses (time zero at first analgesic), cumulative reassessment curves by race/ethnicity largely overlapped (Figure~\ref{fig:postrx}).
The multivariate log-rank test was not statistically significant in the primary trauma pain-score-10 subgroup, suggesting that demographic differences in the primary analysis are not primarily driven by post-treatment reassessment intervals.
This supports interpreting insurance and race associations as operating earlier in the encounter or through pathways other than post-analgesic documentation alone.

\begin{figure}[htbp]
  \centering
  \includegraphics[width=0.85\linewidth]{fig13_post_analgesic_pathway.png}
  \caption{Pathway sensitivity: post-analgesic reassessment by race/ethnicity (time zero = first analgesic).}
  \label{fig:postrx}
\end{figure}
